\subsection*{II.}
% unit: noether.p08.par.009
在证明归约定理以前，先回忆关于线性形式族的一些事实。所谓 \emph{$K$ 上的线性形式族}，是指同一维数的形式系统，它是完备的，意即只要含有 $\theta_1$ 和 $\theta_2$，就总含有 $c_1\theta_1+c_2\theta_2$，其中 $c_1,c_2$ 是给定有理性域 $K$ 中的量，并且 $K$ 包含各 $\theta$ 的系数域。在这些形式中，总有有限个（设为 $\rho$ 个）在 $K$ 上线性无关，而任意 $\rho+1$ 个形式都满足一个系数取自 $K$ 的线性关系；这个形式族关于 $K$ 具有\emph{秩 $\rho$}。\footnote{若去掉同维数的限制，或者去掉 $K$ 包含 $\theta$ 的系数域这一限制，就得到更一般的线性形式族；它们的秩不必是有限的。}我们将使用一个显然的事实：若 $\mathcal L$ 与 $\mathcal T$ 是 $K$ 上两个同秩 $\rho$ 的线性形式族，并且 $\mathcal T$ 是 $\mathcal L$ 的子族，则 $\mathcal L$ 与 $\mathcal T$ 相同（等价）。

% unit: noether.p08.par.010
有了这些，我们转向归约定理。I 中的恒等式
% unit: noether.p08.eq.009
\[
  \theta=\sum PZ
\]
在对两边施加同一个极化过程 $P$ 时，变为 $P\theta$ 的类似恒等式；因此归约定理同时给出了 $\theta$ 的所有极化形式用特殊极化形式 $PZ$ 的和来表示。反过来，这些特殊极化形式当然包含在所有极化形式中，所以归约定理断言这两个线性形式族等价；特别也断言其中每个形式在各个行中的维数均与 $\theta$ 相同的子族等价。为证明这种等价，我们插入由形式 $Z$ 决定的中间形式族。为简单起见，先给出三行二元的情形 $(N=3,n=2)$ 的证明，然后简要指出完全平行的一般证明。

% unit: noether.p08.par.011
于是令 $\theta(x,y,z)$ 在行
% unit: noether.p08.eq.010
\[
  x_1x_2;\quad y_1y_2;\quad z_1z_2;
\]
中分别具有维数 $\alpha,\beta,p$；并且先把 $\theta$ 的系数看作不定元 $a$（或有理数）。我们考虑以下三个线性形式族。

% unit: noether.p08.par.012
1. 形式族 $\mathcal L$，由所有从 $\theta$ 通过极化过程 $P$ 得到、且在各个行 $x,y,z$ 中具有与 $\theta$ 相同维数的形式 $f(x,y,z)$ 构成；特别地，$\mathcal L$ 包含 $\theta$。若把 $f(x,y,z)$ 看作 $x,y,z$ 以及不定元 $a$ 中的形式，则系数域是有理数域 $R$；$K$ 也必须取为 $R$，因为只有对整有理的 $\theta_1,\theta_2$，$c_1P_1+c_2P_2$ 才重新成为一个过程 $P$；设 $\mathcal L$ 关于 $R$ 的秩为 $\rho$。

% unit: noether.p08.par.013
2. 形式族 $\mathcal S$，由所有按
% unit: noether.p08.eq.011
\[
  g(\lambda;x,y)=f(x,y,\lambda_1x+\lambda_2y)
\]
定义的形式 $g(\lambda;x,y)$ 构成，或用极化过程表示为
% unit: noether.p08.eq.012
\[
\begin{aligned}
  g(\lambda;x,y)
  &=\frac1{p!}\left[(\lambda_1x+\lambda_2y)\frac{\partial}{\partial z}\right]^p f(x,y,z) \\
  &=g_0(xy)\lambda_1^p+g_1(x,y)\lambda_1^{p-1}\lambda_2+\cdots+g_p(xy)\lambda_2^p,
\end{aligned}
\]
其中
% unit: noether.p08.eq.013
\[
  i!(p-i)!\,g_i(xy)=
  \left(y\frac{\partial}{\partial z}\right)^i
  \left(x\frac{\partial}{\partial z}\right)^{p-i}f(xyz).
\]
\footnote{如 I 中一样，
\[
 \left(t\frac{\partial}{\partial x}\right)=t_1\frac{\partial}{\partial x_1}+t_2\frac{\partial}{\partial x_2},
\]
因而特别有
\[
 \left[(\lambda_1x+\lambda_2y)\frac{\partial}{\partial z}\right]
 = (\lambda_1x_1+\lambda_2y_1)\frac{\partial}{\partial z_1}
 + (\lambda_1x_2+\lambda_2y_2)\frac{\partial}{\partial z_2},
\]
\[
 \left[z\left(\frac{\partial}{\partial\lambda_1}\frac{\partial}{\partial x}
 +\frac{\partial}{\partial\lambda_2}\frac{\partial}{\partial y}\right)\right]
 =z_1\left(\frac{\partial}{\partial\lambda_1}\frac{\partial}{\partial x_1}
 +\frac{\partial}{\partial\lambda_2}\frac{\partial}{\partial y_1}\right)
 +z_2\left(\frac{\partial}{\partial\lambda_1}\frac{\partial}{\partial x_2}
 +\frac{\partial}{\partial\lambda_2}\frac{\partial}{\partial y_2}\right).
\]}
这样，$g_i(xy)$ 给出恒等式 (1) 中的形式 $Z$。各 $g$ 是 $x,y,\lambda$ 以及不定元 $a$ 中的有理整系数形式；设 $\mathcal S$ 关于 $R$ 的秩为 $\sigma$。

% unit: noether.p08.par.014
3. 形式族 $\mathcal T$，它由 $\mathcal S$ 经逆极化过程得到，正如 $\mathcal S$ 由 $\mathcal L$ 得到一样；$\mathcal T$ 中的形式 $h$ 定义为
% unit: noether.p08.eq.014
\[
\begin{aligned}
 h(x,y,z)
 &=\frac1{p!}\left[z\left(\frac{\partial}{\partial\lambda_1}\frac{\partial}{\partial x}
 +\frac{\partial}{\partial\lambda_2}\frac{\partial}{\partial y}\right)\right]^p g(\lambda;xy)\\
 &=h_0(xyz)+h_1(xyz)+\cdots+h_p(xyz),
\end{aligned}
\]
其中由 2. 得
% unit: noether.p08.eq.015
\[
  h_i(xyz)=
  \left(z\frac{\partial}{\partial x}\right)^{p-i}
  \left(z\frac{\partial}{\partial y}\right)^i g_i(xy).
\]
各个 $h_i$，从而 $h$ 本身，都是形式 $PZ$ 以及它们的和；设 $\mathcal T$ 关于 $R$ 的秩为 $\tau$。

% unit: noether.p08.par.015
由 $\mathcal T$ 的形式 $h(x,y,z)$ 的构造可见，这些形式由 $\theta$ 通过极化过程 $P$ 得到，并且在行 $x,y,z$ 中分别具有与 $\theta$ 相同的维数；因此 $\mathcal T$ 是 $\mathcal L$ 的子族。由上面提到的事实，为证明等价只剩下证明秩相等。在此证明中，不再使用 $f(xyz)$ 的特殊结构，即它们是同一形式 $\theta$ 的极化形式。

% unit: noether.p08.par.016
\emph{a) 为比较 $\mathcal L$ 与 $\mathcal S$ 的秩，使用引理 a)：由 $f(x,y,\lambda_1x+\lambda_2y)=0$ 必然推出 $f(x,y,z)=0$。} 这直接来自三个二元行之间的恒等关系
% unit: noether.p08.eq.016
\[
  (xy)z+(yz)x+(zx)y=0,
  \qquad (xy)=(x_1y_2-x_2y_1),\ldots,
\]
它推出
% unit: noether.p08.eq.017
\[
  f[x,y,(xy)x+(xz)y]=(xy)^p f(x,y,z).
\]
因此，由 $f(x,y,\lambda_1x+\lambda_2y)=0$（关于 $\lambda$ 恒等），通过代入 $\lambda_1=(zy)$、$\lambda_2=(xz)$，并且由于 $(xy)\ne0$，必然得 $f(x,y,z)=0$。

% unit: noether.p08.par.017
由这个引理，$\mathcal S$ 中的形式与 $\mathcal L$ 中的形式之间存在一一对应。事实上，由
% unit: noether.p08.eq.018
\[
  f_1(x,y,\lambda_1x+\lambda_2y)=g(\lambda;xy),\qquad
  f_2(x,y,\lambda_1x+\lambda_2y)=g(\lambda;xy)
\]
必然推出
% unit: noether.p08.eq.019
\[
  f_1(xyz)=f_2(xyz).
\]
此外，$\mathcal S$ 中的每一个关系都保留在 $\mathcal L$ 中唯一对应的形式之间（反之亦然）。因为由
% unit: noether.p08.eq.020
\[
  c_1g^{(1)}(\lambda;xy)+c_2g^{(2)}(\lambda;xy)+\cdots+c_rg^{(r)}(\lambda;xy)=0
\]
通过引理必然得到
% unit: noether.p08.eq.021
\[
  c_1f^{(1)}(x,y,z)+c_2f^{(2)}(x,y,z)+\cdots+c_rf^{(r)}(x,y,z)=0.
\]
\emph{由此证明了 $\mathcal L$ 与 $\mathcal S$ 的秩相等；$\rho=\sigma$。}

% unit: noether.p08.par.018
\emph{b) 为比较 $\mathcal S$ 与 $\mathcal T$ 的秩，使用对应的引理 b)：由 $h(x,y,z)=0$ 必然推出 $g(\lambda;xy)=0$。} 证明时注意，由 3.，$h(x,y,z)=0$ 等价于各个幂积
% unit: noether.p08.eq.022
\[
 \left(\frac{\partial}{\partial\lambda_1}\frac{\partial}{\partial x_1}
 +\frac{\partial}{\partial\lambda_2}\frac{\partial}{\partial y_1}\right)^{p_1}
 \left(\frac{\partial}{\partial\lambda_1}\frac{\partial}{\partial x_2}
 +\frac{\partial}{\partial\lambda_2}\frac{\partial}{\partial y_2}\right)^{p_2}g(\lambda;x,y),
 \qquad p_1+p_2=p.
\]
的消失。进一步微分还得到幂积
% unit: noether.p08.eq.023
\[
\begin{aligned}
&\left(\frac{\partial}{\partial x_1}\right)^{\alpha_1}
 \left(\frac{\partial}{\partial x_2}\right)^{\alpha_2}
 \left(\frac{\partial}{\partial y_1}\right)^{\beta_1}
 \left(\frac{\partial}{\partial y_2}\right)^{\beta_2} \\
&\quad\cdot
 \left(\frac{\partial}{\partial\lambda_1}\frac{\partial}{\partial x_1}
 +\frac{\partial}{\partial\lambda_2}\frac{\partial}{\partial y_1}\right)^{p_1}
 \left(\frac{\partial}{\partial\lambda_1}\frac{\partial}{\partial x_2}
 +\frac{\partial}{\partial\lambda_2}\frac{\partial}{\partial y_2}\right)^{p_2}
 g(\lambda;xy)
\end{aligned}
\]
的消失。因此，这些幂积的任意线性组合都消失，从而每一个形式
% unit: noether.p08.eq.024
\[
  f\left(\frac{\partial}{\partial x},\frac{\partial}{\partial y},
  \frac{\partial}{\partial\lambda_1}\frac{\partial}{\partial x}
  +\frac{\partial}{\partial\lambda_2}\frac{\partial}{\partial y}\right)g(\lambda,xy)
\]
也消失。特别选择 $f$ 使得
% unit: noether.p08.eq.025
\[
  f(x,y,\lambda_1x+\lambda_2y)=g(\lambda;xy),
\]
便还得到
% unit: noether.p08.eq.026
\[
  g\left(\frac{\partial}{\partial x};
    \frac{\partial}{\partial x},
    \frac{\partial}{\partial y}\right)g(\lambda,x,y)=0;
\]
或者，如果
% unit: noether.p08.eq.027
\[
  g(\lambda,xy)=\sum G_i\lambda_1^{\nu_1}\lambda_2^{\nu_2}
  x_1^{\alpha_1}x_2^{\alpha_2}y_1^{\beta_1}y_2^{\beta_2},
\]
其中 $G_i$ 是不定元 $a$ 中的有理整系数线性形式，则
% unit: noether.p08.eq.028
\[
  \sum \nu_1!\nu_2!\alpha_1!\alpha_2!\beta_1!\beta_2!\,G_i^2=0,
\]
于是 $g(\lambda;xy)=0$。

% unit: noether.p08.par.019
由引理 b)，完全如 a) 中那样推出：\emph{$\mathcal S$ 与 $\mathcal T$ 的秩相等；$\sigma=\tau$。}

% unit: noether.p08.par.020
因此由 a) 和 b) 得 $\rho=\tau$；于是，由于 $\mathcal T$ 是 $\mathcal L$ 的子族，这两个线性形式族的等价性得到证明。于是 $\mathcal L$ 的每一个形式，特别是 $\theta$，都可表示为 $\rho$ 个线性无关形式 $h(x,y,z)$ 的有理整系数线性组合：
% unit: noether.p08.eq.029
\[
  \theta=\sum c_i h^{(i)}(x,y,z)=\sum PZ.
\]
这个对 $\theta$ 的不定系数 $a$ 推出的恒等式，在把不定元 $a$ 替换为任意有理性域中的量时仍然成立；因此，\emph{三行二元情形下的归约定理已一般地得到证明。}

% unit: noether.p08.par.021
对于 $N$ 个、每个含 $n$ 个变量的行，一般证明完全平行。设 $\theta(A)$ 在行 $A_i$ 中具有维数 $\alpha_i$。我们再次考虑三个线性形式族：

% unit: noether.p08.par.022
1. 形式族 $\mathcal L$，由所有从 $\theta(A)$ 通过极化过程 $P$ 得到、并在每个单独行 $A_i$ 中具有与 $\theta(A)$ 相同维数的形式 $f(A)$ 构成；

% unit: noether.p08.par.023
2. 形式族 $\mathcal S$，由所有通过代入
% unit: noether.p08.eq.030
\[
\begin{aligned}
 A_{n+1}&=\lambda_{n+1}^{(1)}A_1+\lambda_{n+1}^{(2)}A_2+
        \cdots+\lambda_{n+1}^{(n)}A_n,\\
 &\vdotswithin{=}\\
 A_N&=\lambda_N^{(1)}A_1+\lambda_N^{(2)}A_2+
        \cdots+\lambda_N^{(n)}A_n;
\end{aligned}
\]
从 $f(A)$ 得到的形式 $g(\lambda;A_1\cdots A_n)$ 构成；在此过程中，$g(\lambda;A_1\cdots A_n)$ 中各个 $\lambda$ 幂积的系数给出恒等式 (1) 的形式 $Z$；

% unit: noether.p08.par.024
3. 形式族 $\mathcal T$，由所有按
% unit: noether.p08.eq.031
\[
\begin{aligned}
 h(A)=&\left[A_{n+1}\left(
 \frac{\partial}{\partial\lambda_{n+1}^{(1)}}\frac{\partial}{\partial A_1}
 +\cdots+
 \frac{\partial}{\partial\lambda_{n+1}^{(n)}}\frac{\partial}{\partial A_n}\right)\right]^{\alpha_{n+1}} \\
&\cdots
\left[A_N\left(
 \frac{\partial}{\partial\lambda_N^{(1)}}\frac{\partial}{\partial A_1}
 +\cdots+
 \frac{\partial}{\partial\lambda_N^{(n)}}\frac{\partial}{\partial A_n}\right)\right]^{\alpha_N}
 g(\lambda;A_1\cdots A_n)
\end{aligned}
\]
定义的形式 $h(A)$ 构成。这里 $h(A)$ 是形式 $PZ$ 的和。

% unit: noether.p08.par.025
由于任意 $n+1$ 个行之间存在恒等关系，引理 a) 仍然成立；引理 b) 同样仍然成立，因为变量数并未进入它的证明。因此 $\mathcal L$ 与 $\mathcal T$ 具有相同的秩，并且由于 $\mathcal T$ 是 $\mathcal L$ 的子族，它们相同。所以恒等式
% unit: noether.p08.eq.032
\[
  \theta=\sum PZ
\]
对不定系数成立，从而也对任意有理性域中的量作为系数成立；\emph{归约定理由此在一般情形下得到证明。}

% unit: noether.p08.sec.003
